% Supplement: rounding-inclusive computation.
% Source: draft/bie-enclosure/a-reliable.tex (clean version), Sections C.1--C.6.
\section{Rounding-inclusive computation}\label{sm:reliable}
This section describes how the inputs in \cref{tab:inputs} are computed. All
bounds refer to the same field and the same $\widehat k$.

\subsection{Principles}\label{sm:principles}
\emph{Saved data are exact numbers.} The candidate $\widehat k$, density samples,
proxy coefficients and propagation matrices are stored as binary64 numbers and
treated as exact binary rationals; the field of \cref{sec:field} built from them
is the field to which \cref{thm:distance} is applied. Geometry parameters
determine the target operator; the target geometry is defined by the saved
binary64 parameters (\cref{sec:results}). The defects and the mass bound are
computed for this geometry; the gap test and the trace constants are computed for
parameter balls that contain both the binary64 and the decimal values. The
number $\pi$, the phase $q$, the wavenumbers $\widehat k\rho^{1/2}$, the Rayleigh
exponents $\gamma_m$ and special functions are evaluated from exact formulas.
Rounding in the stored field data therefore needs no estimate: it changes which
field is tested, and the theorem holds for every fixed field.

\emph{Floating point proposes, ball arithmetic decides.} Floating point is used
only for proposals: the candidate, fit coefficients, panel splits, trial vectors
and subspaces. Every decision (error bound, sign, definiteness, inequality) is
made in real or complex ball arithmetic \cite{Johansson2017}, where a ball
$[m\pm r]$ contains the exact value and every operation rounds outward. The
accuracy of the floating-point results affects only the tightness of the bounds.

\emph{Sum with signs first.} An interface jump is a difference of two quantities
of order one. Approximations of the parts are therefore added with their signs,
and only the error radii are added in absolute value; for Taylor or Chebyshev
models the coefficients are combined before the polynomial is integrated. Printed
decimal endpoints are rounded outward.

\subsection{Interface defects}\label{sm:reliable-eps}
Every interface is split into panels, and on each panel a polynomial $P_I$ and an
error $e_I\ge\sup_I\|f-P_I\|$ for the weighted jump vector $f$ of \eqref{eq:panel}
are computed. In this section the curve parameter has period $2\pi$, panels are
$I=[c-h,c+h]$, and vectors and matrices carry the Euclidean and Frobenius norms,
so the bounds hold verbatim for scalar jumps and for matrices mapping lead states
to jumps.

\begin{lemma}[Chebyshev interpolation with node errors]\label{sm:lem-cheb}
Let $f$ be holomorphic near the closed Bernstein ellipse $E_\varrho$ ($\varrho>1$)
with foci at the ends of $I$, and $\sup_{E_\varrho}\|f\|\le M_I$. Let
$t_l=c+h\cos\frac{(2l+1)\pi}{2L}$, $0\le l<L$, and let $v_l$ satisfy
$\|v_l-f(t_l)\|\le e_{\mathrm{node}}$. Compute the coefficients of the degree $L-1$
interpolant of the $v_l$ in ball arithmetic, take their midpoints as $P_I$, and let
$e_{\mathrm{conv}}$ be the sum of the radii. Then
\begin{equation}\label{sm:eq-cheb}
  \sup_{t\in I}\|f(t)-P_I(t)\|\le\frac{4M_I\varrho^{1-L}}{\varrho-1}
  +\sqrt{2L-1}\,e_{\mathrm{node}}+e_{\mathrm{conv}} .
\end{equation}
\end{lemma}

\begin{proof}
With $t=c+h(w+w^{-1})/2$, the Cauchy formula on $|w|=\varrho$ gives Chebyshev
coefficients $\|a_k\|\le2M_I\varrho^{-k}$, $k\ge1$. At the $L$ roots, each $T_k$
with $k\ge L$ coincides with $\pm T_j$, $j<L$, or with zero, so for exact node
values the interpolation error is at most $2\sum_{k\ge L}2M_I\varrho^{-k}$. The
discrete cosine coefficients $d_j$ of the node errors satisfy
$\|d_0\|^2+\frac12\sum_{j\ge1}\|d_j\|^2=L^{-1}\sum_l\|e_l\|^2\le e_{\mathrm{node}}^2$,
and $|T_j|\le1$ with Cauchy--Schwarz bounds their interpolant by
$(2L-1)^{1/2}e_{\mathrm{node}}$. The exact coefficients lie in the computed balls
and the basis functions are bounded by one on $I$, which gives $e_{\mathrm{conv}}$.
\end{proof}

Fixing the node midpoints before interpolation propagates node errors once,
through $\sqrt{2L-1}$; transforming balls directly would amplify the radii. Keeping
the signs of all parts in $P_I$ retains the cancellation of about twelve digits in
the material jumps. Node values have three parts: the self-interaction of the
curve's own layer potentials (Kress splitting, \cref{sm:lem-self,sm:lem-kress});
the smooth part, namely the two near images and the proxies, by regular expansions
(\cref{sm:lem-regular}); and local terms, namely continuous density values and
background waves (\cref{sm:lem-local}).

\begin{lemma}[regular Helmholtz expansion]\label{sm:lem-regular}
For a real center $\mathbf x_\nu$ let $z=(x-x_\nu)+\ii(y-y_\nu)$ and
$w=(x-x_\nu)-\ii(y-y_\nu)$, complexified as independent variables,
$\lambda=\kappa^2/4$, and $\Phi_m(\xi)=\sum_{j\ge0}m!\,\xi^j/(j!\,(m+j)!)$. Let $Y$ be
holomorphic near the bidisk $|z|,|w|\le R$, with each component solving
$4\partial_z\partial_wY+\kappa^2Y=0$ and $\|Y\|\le B$. Then
\[
  Y=a_0\Phi_0(-\lambda zw)+\sum_{m\ge1}\bigl[a_m^+(z/R)^m+a_m^-(w/R)^m\bigr]\Phi_m(-\lambda zw),
\]
with $\|a_m^\pm\|\le B$. If the series is cut at order $M$ and the coefficients are
replaced by computed values with errors $\epsilon_0,\epsilon_m^\pm$, then on real
targets with $|z|=|w|\le a<R$, $\theta=a/R$, the error is at most
\[
  \ee^{\lambda a^2}\Bigl[\frac{2B\theta^{M+1}}{1-\theta}+\epsilon_0
  +\sum_{m=1}^M\theta^m(\epsilon_m^++\epsilon_m^-)\Bigr].
\]
Taking $Y=(V,\partial_xV,\partial_yV)$ bounds values and gradients together.
\end{lemma}

\begin{proof}
Writing $Y=\sum C_{ab}z^aw^b$, the equation gives
$C_{a+1,b+1}=-\lambda C_{ab}/((a+1)(b+1))$, so each diagonal is determined by a pure
derivative; summing along diagonals gives the expansion, and the Cauchy formula
gives $\|a_m^\pm\|\le B$. Since $m!/(m+j)!\le1$, $|\Phi_m(-\lambda zw)|\le
\ee^{\lambda a^2}$; sum the two geometric tails.
\end{proof}

The proxies need one expansion about the origin with $R=1.1$. The near images use
$R=0.6$ about the origin for material targets and the panel midpoints for walls and
seams; their coefficients are source integrals, computed from source-panel moments.
On source panels $S_s=[c_s-h_s,c_s+h_s]$ let the kernel (including the source
Jacobian and the Müller signs) be holomorphic and bounded by $M_{\nu s}$ on the
bidisk times $|s-c_s|\le R_s$, $q_s=h_s/R_s<1$, and $D_s\ge\int_{S_s}\|\eta\|$.
Truncating the kernel at degree $n_s$ in $s$ and using the normalized moments
$\int_{S_s}((s-c_s)/h_s)^j\eta\,\dd s$ with errors $d_{s,j}$ gives an error at most
\begin{equation}\label{sm:eq-moment}
  E_\nu=\sum_sM_{\nu s}\Bigl[D_s\frac{q_s^{n_s+1}}{1-q_s}+\sum_{j\le n_s}q_s^jd_{s,j}\Bigr]
\end{equation}
on the bidisk, with $B=\sum_sM_{\nu s}D_s$. The truncated expansion still solves the
Helmholtz equation term by term, so $E_\nu(1+2\sum_{m\le M}\theta^m)$ is added to the
error above. The series of $\Phi_m$ is truncated with a geometric bound on the
remainder.

\begin{lemma}[envelope of the fundamental solution]\label{sm:lem-green}
Let $G(q)=\frac\ii4H_0^{(1)}(\kappa\sqrt q)$ as a function of the squared distance
$q$, $b_0=\kappa^2/4$, $C_0=\frac14+(|\log(\kappa/2)|+\gamma_{\mathrm E})/(2\pi)$. On a
complex domain with $q_-\le|q|\le Q$ and a branch with $|\log q|\le L_q$, let
$\mathcal G_j=b_0^j\ee^{b_0Q}(L_q/(4\pi)+C_0+(j+b_0Q)/(2\pi))$. Then
\[
  |G|\le\mathcal G_0=:g_0,\quad
  |G'|\le\mathcal G_1+\frac{\ee^{b_0Q}}{4\pi q_-}=:g_1,\quad
  |G''|\le\mathcal G_2+\frac{\ee^{b_0Q}}{4\pi}\Bigl(\frac{2b_0}{q_-}+\frac1{q_-^2}\Bigr)=:g_2 ,
\]
and for displacements $\|\boldsymbol d\|\le D$, $\|\nabla G\|\le2Dg_1$ and
$\|\nabla^2G\|\le2g_1+4D^2g_2$.
\end{lemma}

\begin{proof}
From the series of $Y_0$,
\[
  G(q)=-\frac{\log q}{4\pi}J(q)+\Bigl[\frac\ii4-\frac{\log(\kappa/2)+\gamma_{\mathrm E}}{2\pi}\Bigr]J(q)
  +\frac1{2\pi}\sum_{m\ge1}\frac{(-1)^mH_m(b_0q)^m}{(m!)^2},
\]
with $J(q)=\sum_{m\ge0}(-b_0q)^m/(m!)^2$ and harmonic numbers $H_m\le m$. Bound term by term;
$\nabla G=2G'\boldsymbol d$ and $\nabla^2G=2G'I+4G''\boldsymbol d\boldsymbol d^{\top}$.
\end{proof}

For source disks at distance $d_-\le d\le d_+$ from the target center, source
displacement $\delta$ and radius $R$ with $R+\delta<d_-$, take
$q_-=(d_--R-\delta)^2$, $Q=(d_++R+\delta)^2$ and
$L_q=\max\{|\log d_-^2|,|\log d_+^2|\}-2\log(1-(R+\delta)/d_-)$; the logarithm is the
real value at the center plus two convergent $\log(1+u)$ series, on the branch of
the real kernel. If the source speed is at most $V_s$ and the normal at most $N_s$,
then with $D=d_++R+\delta$ a group of point sources with coefficients $c_j$ gives
\begin{equation}\label{sm:eq-kernel}
  M_{\nu s}\le V_s\sum_j|c_j|\bigl[g_0+2(N_s+1)Dg_1+N_s(2g_1+4D^2g_2)\bigr].
\end{equation}

\begin{lemma}[local terms]\label{sm:lem-local}
A background wave $b\exp(\ii\phi(t))$ with $\phi=\pm\gamma(x(t)-x_a)+\beta y(t)$ has
value and gradient bounded on the complex target domain by
\[
  |b|\exp\bigl(\sup\Re(\ii\phi)\bigr)\bigl(1+|\gamma|^2+\beta^2\bigr)^{1/2},
\]
and a Fourier density
$\sum\eta_m\ee^{\ii mt}$ is bounded on a target disk of radius $R_t$ by
$\sum\|\eta_m\|\ee^{|m|R_t}$. Multiplying by the bound of the weight and summing gives
$M_{\mathrm{loc}}$. For Taylor truncation at degree $p_{\mathrm{loc}}$ and $h<R_t$,
\[
  e_{\mathrm{loc}}\le M_{\mathrm{loc}}\frac{(h/R_t)^{p_{\mathrm{loc}}+1}}{1-h/R_t}
  +e_{\mathrm{loc,fin}},
\]
where $e_{\mathrm{loc,fin}}$ is the radius of the finite coefficients.
\end{lemma}

The proof is the Cauchy coefficient bound $M_{\mathrm{loc}}R_t^{-k}$. On a center
panel, $e_I=e_{\mathrm{self}}+e_{\mathrm{loc}}+e_{\mathrm{sm}}+e_{\mathrm{comb}}$, where
$e_{\mathrm{sm}}$ is \eqref{sm:eq-cheb} for the smooth part ($L=48$, $\varrho=4.5$)
with expansion tails and node radii in $e_{\mathrm{node}}$, and $e_{\mathrm{comb}}$ is
the radius dropped when the three coefficient vectors are added and rounded to
midpoints. The coefficients of $P_I$ are then exact binary64 numbers, and
$V_I=\int_I\|P_I\|^2$ is computed from exact monomial integrals in ball arithmetic
and checked by Gauss--Legendre quadrature. The weight $|\boldsymbol r'|^{1/2}$ and
the factor $1/|\boldsymbol r'|$ in the normal derivative are continued as $s^{1/4}$
and $s^{-1/2}$ of $s(t)=x'(t)^2+y'(t)^2$ on the principal branch, after checking
$\Re s\ge\underline s>0$ on every complex strip used.

\emph{Self-interaction.} In the normal jump on a rod, the two hypersingular parts
cancel analytically in the Müller difference kernel. The remaining self term is
\begin{equation}\label{sm:eq-self}
  F(t)=\frac1{2\pi}\int_0^{2\pi}\Bigl\{A(t,s)\log\Bigl(4\sin^2\frac{t-s}2\Bigr)+B(t,s)\Bigr\}\,\dd s,
\end{equation}
where $A,B$ contain the density, Jacobians and weights. Let $A,B$ be holomorphic for
targets in $T$ and sources in $|\Im s|\le a$, with $\|A\|\le H_A$, $\|B\|\le H_B$;
before use we check in complex balls that the squared speed and the chord quotient
in the logarithm have positive real part, so all branches agree with the real
kernel. For complex targets \eqref{sm:eq-self} means its Fourier continuation
$F(t)=-\sum_{m\ne0}\widehat A_m(t)\ee^{\ii mt}/|m|+\widehat B_0(t)$, and shifting the
contour gives $\|\widehat A_m(t)\|\le H_A\ee^{-a|m|}$.

\begin{lemma}[target envelope]\label{sm:lem-self}
If $T$ contains the disk $|t-c|\le R_t<a$, then $\|F\|\le2H_A/(\ee^{a-R_t}-1)+H_B$ there.
\end{lemma}

This follows from $|\ee^{\ii mt}|\le\ee^{|m|R_t}$. If $h(\varrho+\varrho^{-1})/2<R_t$,
it serves as $M_I$ in \cref{sm:lem-cheb}. Node values use Kress quadrature with
$N=2n$ equispaced sources: the trigonometric interpolant of $A(t,\cdot)$ (Nyquist
mode split evenly) is convolved exactly with the logarithm, and $B$ uses the
trapezoidal rule.

\begin{lemma}[Kress error at real targets]\label{sm:lem-kress}
Let $k_N(m)$ be the absolute value of the representative of $m$ modulo $N$ in
$[-n,n]$, $w_N(m)=1/k_N(m)$ if $k_N(m)>0$ and $0$ otherwise. Then the Kress error
is at most
\[
  e_{\mathrm{Kress}}\le2H_A\Bigl[\frac{\ee^{-an}}{n(1-\ee^{-a})}
  +\frac{\ee^{-an}}{1-\ee^{-aN}}\sum_{j=0}^{N-1}\ee^{-aj}w_N(n+j)\Bigr]
  +\frac{2H_B}{\ee^{aN}-1}.
\]
\end{lemma}

\begin{proof}
Modes $|m|<n$ are reproduced exactly. A mode $|m|\ge n$ contributes at most
$\|\widehat A_m\|/|m|$ to the exact convolution and, aliased to $k_N(m)$, at most
$w_N(m)\|\widehat A_m\|$ to the quadrature. Sum over $m\ge n$ on both sides, using
$1/m\le1/n$ in the first part and the period $N$ of $w_N$ in the second. The
trapezoidal error is from \cite[Thm.~3.2]{TrefethenWeideman2014}.
\end{proof}

With node balls of maximal radius $e_{\mathrm{rad}}$ at $L$ roots,
\cref{sm:lem-cheb} gives
$e_{\mathrm{self}}\le4M_{\mathrm{self}}\varrho^{1-L}/(\varrho-1)+\sqrt{2L-1}(e_{\mathrm{Kress}}
+e_{\mathrm{rad}})+e_{\mathrm{conv}}$. On the center curve $N=1152$, the envelope uses
$R_t=0.04$, $a=0.0625$, the Kress error $a=0.08$, and $L$ is increased from $40$
until the interpolation remainder is below $2\times10^{-17}$ (basis change at $192$
bits). Lead curves use $R_t=0.12$, $a=0.16$.

\emph{Lead cells.} On a lead panel the jump depends linearly on the state,
$f_{I,j}=F_I(t)v_j$ with $v_j=\mathsf P^jc$. For $j<N_{\mathrm{cell}}$ we build
$P_{I,j}$ and $e_{I,j}$ for each actual state and use
$\sum_j\|f_{I,j}\|^2_{L^2(I)}\le\sum_j(V_{I,j}^{1/2}+|I|^{1/2}e_{I,j})^2$. The tail
\eqref{eq:tail} uses $\Theta_\pm=B_\pm^2$ with analytic envelopes per group and side
(for unit states)
\begin{gather*}
  B_{\mathrm{mat}}=\sqrt{2\pi}\bigl[(S_{\mathrm{self}}+S_{\mathrm{loc}}+wS_{\mathrm{near}})
  \|\mathsf G\|+wS_{\mathrm{pt}}\bigr],\\
  B_{\mathrm{wall}}=(1+r)S_{\mathrm{face}},\qquad B_{\mathrm{seam}}=2S_{\mathrm{face}},
\end{gather*}
$S_{\mathrm{face}}=(S_{\mathrm{ord}}+S_{\mathrm{inc}})\|\mathsf G\|+S_{\mathrm{pt}}$. Here
$\mathsf G$ maps states to density coefficients; $S_{\mathrm{near}}$, $S_{\mathrm{ord}}$
and $S_{\mathrm{pt}}$ bound the near-image potentials at material targets, the
ordinary potentials (three images) at walls and seams, and the $480$ point sources,
all through \eqref{sm:eq-kernel}; $S_{\mathrm{self}}$ is the maximum of
\cref{sm:lem-self} over the $64$ panels ($a=0.16$, $h=\pi/64$); $S_{\mathrm{loc}}$
and $S_{\mathrm{inc}}$ bound the local and incident Rayleigh terms, the latter by
$(\sum_{|m|\le48}2(1+|\gamma_m|^2+\beta_m^2)/d)^{1/2}$; and
$w=\overline v^{1/2}\max\{1,\overline v/\underline s^{1/2}\}$ comes from the weight
($w<1$ is verified). The factor $\sqrt{2\pi}$ turns pointwise bounds into $L^2$
bounds; a wall involves $v_j$ and $\mathsf Pv_j$, a seam two traces. These envelopes
are much cruder than the actual defects (\cref{sm:tab-tails}) but enter only the
tail. For each group and side we take the smallest $N_{\mathrm{cell}}$ with tail at
most $10^{-6}$ times the lower bound of the finite part, and check the decay with
$N_{\mathrm{cell}}+10$.

\emph{Assembly.} Walls between the center and the first lead cell combine the
polynomials of both sides; interior walls use two consecutive states; each seam pair
combines top, bottom and phase once. Then
\begin{equation}\label{sm:eq-eps}
  \overline\varepsilon_g^{\,2}=\sum_{(I,j)\in\mathscr P_g}\bigl(V_{I,j}^{1/2}+|I|^{1/2}e_{I,j}\bigr)^2
  +\overline T_g ,
\end{equation}
where $\mathscr P_g$ contains all (panel, state) pairs of the group (one state for
center panels) and $\overline T_g$ is the sum of the two tail bounds; the square
root is rounded up to binary64.

\subsection{Mass lower bound}\label{sm:reliable-mass}
Take disjoint closed disks, each inside one homogeneous piece. For each disk we
check in ball arithmetic: (i) its distance to all interfaces exceeds $r$, and the
center is on the right side (along a path); (ii) the disks are pairwise disjoint;
(iii) a value $\widetilde u(\mathbf c)$ with $|\widehat u(\mathbf c)-\widetilde
u(\mathbf c)|\le e_{\mathbf c}$, by Gauss quadrature with an analytic error bound
(the integrands are smooth because $\mathbf c$ is far from the curves). By
\eqref{eq:disk},
\[
  \underline U^{\,2}=\sum\rho_b\pi r^2\bigl[J_0(\kappa_br)^2+J_1(\kappa_br)^2\bigr]
  \max\{0,|\widetilde u(\mathbf c)|-e_{\mathbf c}\}^2\le U^2 ,
\]
with lower ends of the Bessel balls. Two or three digits suffice, since
$\overline Q$ is eleven orders of magnitude smaller than $\underline U$.

\subsection{Trace constants}\label{sm:reliable-cg}
The conditions of \cref{prop:trace} are checked from the analytic curve
parametrizations, not from node samples.

\emph{Injectivity.} If a closed curve has winding number one about the origin,
$|\boldsymbol r'|>0$, $\boldsymbol r\times\boldsymbol r'>0$, the signed distance
$\boldsymbol r\times\boldsymbol r'/|\boldsymbol r'|$ from the origin to every tangent
line is at least $H>\delta$, and $\delta\kappa_*<1$, then $\boldsymbol r+t\nu$ is
injective for $|t|\le\delta$: the polar angle is strictly monotone along each offset
curve, and at fixed angle the distance to the origin is strictly monotone in $t$.

\emph{Curvature.} On parameter intervals $[c-h,c+h]$ we bound
$\sup|\kappa|\le|\kappa(c)|+h\sup|\kappa'|$ in ball arithmetic and bisect failing
intervals up to sixteen levels; the leaves cover the full period. This also gives a
lower bound for the speed.

\emph{Overlap.} If a curve lies in $|x|\le X$, $|y|\le Y$, the collars of all its
lattice translates $(jL,kd)$ are disjoint when $L>2X+2\delta$ and $d>2Y+2\delta$, so
$N_g=1$; the half guides are subfamilies of the full lattice. Walls need $2\delta<L$,
seams $2\delta<d$. Seams form their own group, so their collars may overlap the wall
collars without increasing $N_g$.

\emph{Supremum.} The numerator of the derivative of
$F(\lambda)=p_g(\lambda)/(\lambda+\mu)^2$ is linear, so $\sup_{\lambda\ge\lambda_0}F$ is
the maximum of $F(\lambda_0)$, $c_g$ and $F(\lambda_*)$ with
$\lambda_*=(2a_g-b_g\mu)/(2c_g\mu-b_g)$ if $\lambda_*\ge\lambda_0$. If ball arithmetic
cannot exclude $2c_g\mu=b_g$, $F$ is monotone and the supremum is
$\max\{F(\lambda_0),c_g\}$. The parameters $\delta_g,\eta_g$ need not be optimal.

\subsection{Matrix tests for the common gap}\label{sm:band-tests}
\emph{Quadrature for the matrix entries.} With wave vectors $\xi_i$ of $\psi_i$,
$(B_{\mathcal I})_{ij}=b_{\xi_j-\xi_i}/(L_sd)$, where
$b_\xi=\int_{\Omega_s}\ee^{\ii\xi\cdot\mathbf z}\,\dd\mathbf z$; $b_0=|\Omega_s|$ and, for
$\xi\ne0$, by the divergence theorem,
\[
  b_\xi=-\frac{\ii}{|\xi|^2}\int_0^{2\pi}\ee^{\ii\xi\cdot\boldsymbol c(t)}
  \bigl(\xi_xc_y'(t)-\xi_yc_x'(t)\bigr)\,\dd t .
\]
For the trigonometric curves used here the integrand is entire and $2\pi$-periodic;
if it is bounded by $\mathcal M$ on $|\Im t|\le\varsigma$, the $J$-point trapezoidal
error is at most $4\pi\mathcal M/(\ee^{\varsigma J}-1)$. (A general $C^2$ curve needs
a separate analytic continuation.) An entrywise error $\epsilon_b$ gives
$\|B-\widetilde B\|\le e_B:=|\mathcal I|\epsilon_b$, hence
$\|M-\widetilde M\|\le|\delta_s|e_B$ and $\|G-\widetilde G\|\le\delta_s^2e_B(3+e_B)$. The
lower test uses $\widetilde T+\lambda|\delta_s|e_BI$, the upper test
$\widetilde T^--[\lambda|\delta_s|e_B+\lambda^2\delta_s^2e_B(3+e_B)/c_{\mathcal I}]I$.

\emph{Signs of finite matrices.} For $n_s=1$ two decisions are needed. Lower test: a
binary64 trial vector $x$ with $x^*Hx<0$ in ball arithmetic gives $n(H)\ge1$. Upper
test: take all floating-point eigenvectors of $H$ except the lowest as the columns
of an exact matrix $W$, compute $C=W^*HW$ and $D=W^*W$ in ball arithmetic, and verify
$C\succeq gI$, $0\prec D\preceq dI$ with $g>0$ by Hermitian row-sum bounds. By
Courant--Fischer, $\lambda_2(H)\ge\min_z z^*Cz/z^*Dz\ge g/d>0$, so $n(H)\le1$. The
eigenvectors only select the subspace. As an independent check, ball LDL$^*$
factorizations on some panels have pivot balls that exclude zero. Panel endpoints
are integer multiples of $\pi/(L_sN)$, and integer comparisons check that they cover
$[0,\pi/L_s]$ (with reflection symmetry) or $[-\pi/L_s,\pi/L_s]$.

\subsection{Assembly}\label{sm:assembly}
$\mu$ is the square of the binary64 $\widehat k$ in ball arithmetic. With the published
decimal upper bounds $\overline C_g$ and the rounded-up binary64 values
$\overline\varepsilon_g$ read as exact numbers, we compute
$\overline Q=\sum_g\overline C_g\overline\varepsilon_g$, $\underline D=\underline U-
\overline Q$ (which must be positive) and $R_{\mathrm{spec}}=2\mu\overline Q/\underline D$,
take upper ends, and round the interval endpoints outward. Condition
\eqref{eq:gap-condition} is checked as $\mu-R_{\mathrm{spec}}-\overline a_s>0$ and
$\underline b_s-\mu-R_{\mathrm{spec}}>0$. The assembly is repeated independently in
exact rational arithmetic, and the printed radii and endpoints must agree.
